\section{Experimental details}
\label{app:details}

All experiments run on a CPU with two cores and 7\,GB of memory, using PyTorch and Hugging Face Transformers with the public checkpoints of GPT-2 small \citep{radford2019language} and Pythia \citep{biderman2023pythia}. The repository contains one script per experiment, the exact command lines, the raw JSON outputs and the plotting scripts. The 40 prompts used in Experiments~1 and~3 are short neutral sentence openings written for this study (for example, ``The weather today is cold and cloudy, with a chance of''), and the full list is in the repository.

\subsection{Experiment 1: finite-time Lyapunov exponents}
\label{app:exp1}

For each prompt we take the first $T$ tokens (prompts shorter than $T$ are repeated) and run the model once to record the residual stream $x^0,\dots,x^L$ before the final layer normalization. We then call each block of the Hugging Face implementation directly as a function $F_\ell:\mathbb{R}^{T\times d}\to\mathbb{R}^{T\times d}$, with rotary embeddings and an explicit causal mask, and check that chaining these calls reproduces the model's hidden states (the maximum relative deviation was zero in float32 for every checkpoint). The largest singular value of the Jacobian $DF_\ell(x^\ell)$ is estimated by power iteration on $DF_\ell^{\top}DF_\ell$, using forward-mode (JVP) and reverse-mode (VJP) products from \texttt{torch.func}. The same procedure applied to the composition $F_{L-1}\circ\cdots\circ F_0$ gives the depth-averaged exponent $\frac{1}{L}\log\sigma_{\max}(DF_{0\to L})$. We use eight prompts, contexts of 16 tokens and 10 power iterations per estimate, for eleven Pythia-70M checkpoints (steps 0, 128, 512, 1k, 2k, 4k, 8k, 16k, 32k, 64k and 143k) and six Pythia-160M checkpoints (steps 0, 512, 2k, 8k, 32k and 143k).

\subsection{Experiment 2: bifurcation diagram}
\label{app:exp2}

\emph{Data.} A first-order Markov chain over $V=12$ symbols, with transition rows drawn from a Dirichlet distribution with concentration $0.3$ (seed 0), generates $N=64$ sequences of length 13, and the model predicts the next symbol at each of the 12 positions. \emph{Model.} One pre-LN transformer block with causal attention (two heads, width 32, MLP width 64), learned positional embeddings and a linear readout, in float64. Attention is written out explicitly so that Hessian-vector products work on a CPU. \emph{Training.} Full-batch gradient descent on the mean cross-entropy, from the same initialization for every step size within a seed. Seeds 0, 1 and 2 change both the Markov chain and the initialization. Each run lasts $\max(3000,\,600/\eta)$ steps, capped at 30{,}000, so that all runs cover a comparable amount of gradient-flow time. A run is marked as diverged if the loss exceeds 50 or becomes non-finite. \emph{Measurements.} The last 400 losses of each run, and the top Hessian eigenvalue at the final iterate from 40 steps of power iteration on Hessian-vector products. The step sizes are 56 values spaced geometrically between $0.02$ and $4$.

\subsection{Experiment 3: escape from repetition loops}
\label{app:exp3}

\emph{Loop detection.} From each prompt we decode greedily for 320 tokens and store the full logit vector at each step. A loop of period $k\le64$ is detected when the last $3k$ tokens are $k$-periodic, and its onset is the earliest index from which the sequence stays $k$-periodic. \emph{Stay probabilities.} Along the loop the context is deterministic, so the probability \eqref{eq:stay} of producing the loop token at temperature $T$ follows from the stored logits, for 30 temperatures spaced geometrically between $0.15$ and $2$. From these we compute, for each cycle $c$, the barrier $\Delta_c$ (smallest gap between the loop token's logit and any other logit during the cycle) and the escape probability $r_c(T)=1-\prod_{t\in c}p_t(T)$. \emph{Arrhenius fit.} For the first cycle after two completed cycles of burn-in, we fit $\log r$ against $1/T$ by least squares over the temperatures where $10^{-9}<r<10^{-2}$, and compare the negative slope with $\Delta_{\min}$ of that cycle. \emph{Monte Carlo check.} For the first four loops we sample 128 continuations at $T\in\{0.6,1.0\}$ (nearest grid temperature) from the context after burn-in, record the first step whose token differs from the loop, and compare the fraction still in the loop after 1, 2 and 4 cycles with the exact survival probability. \emph{Free sampling.} For each of seven temperatures we sample 200 tokens from every prompt with plain temperature sampling (no top-$k$ or top-$p$) and record whether the output ends in a loop with at least three repetitions.
